 %-----------EXPERIENCE-----------%
\section{Experiences}


\resumeSubHeadingListStart
    \resumeSubheading
    {HUINNO}{Dec 2025 -- Present}
    {Data Engineer}{서울}
    \resumeItemListStart
        \resumeItem{\textbf{사내 MLOps 파이프라인 전체를 처음부터 설계·구축·운영}(수집 → 레이크하우스 → 변환 → 계보 → 분산 학습 → 실험·모델 레지스트리). \textbf{플랫폼 2인 팀에서 클러스터 구성 전 구간을 단독 수행} — 네트워크·스토리지·카탈로그·계보·오케스트레이션·분산 학습·신원까지 12개 Pulumi 스택 전부를 직접 만짐. 6개월간 \textbf{10노드 · 상시 pod 235개 · 사내 서비스 23종}을 프로덕션 운영하며 머지 PR 678건 · 커밋 1,204건을 냈고, 베어메탈 \textbf{Kubernetes}를 \textbf{Pulumi}로 전면 코드화해 전 구간을 재현 가능한 배포로 전환.}
        \resumeItem{\textbf{ETL 파이프라인 전 갈래를 직접 설계·구현}(2026.03--07, 5개월 · 별도 저장소, 커밋 1,599건). 제품 공용 프레임워크와 CAST · MIMIC-IV · 운영 유지보수 파이프라인을 직접 작성했고, 나머지 제품도 이 프레임워크 위에서 확장되도록 설계 — \textbf{도메인 ETL은 연구자가 얹고 실행·계보·검증·재시도는 플랫폼이 책임지는} 분담 구조를 정착.}
        \resumeItem{\textbf{[데이터 · 2026.03--07] 데이터셋 준비의 1인 의존을 제거.} 기존에는 연구자가 데이터 과학자 한 명에게 요청 → 하드코딩된 스크립트를 SSH로 \textbf{수 시간} 실행 → 파일로 전달하는 구조라 담당자 부재 시 작업이 멈췄고 재현·계보·버전이 모두 없었다. \textbf{Airflow 3.x + dbt(Cosmos) + Apache Iceberg}로 Bronze/Silver/Gold 변환을 \textbf{스케줄 자동 실행}으로 바꾸고 \textbf{DataHub}에 테이블 단위 계보를 자동 기록해, 요청·대기 없이 연구자가 직접 조회하는 흐름으로 전환.}
        \resumeItem{\textbf{[학습 · 2026.08--현재] 단일 GPU 수동 실행을 대체할 분산 학습 기반을 구축.} \textbf{KubeRay}로 사양이 제각각인 \textbf{GPU 10장}을 단일 풀로 묶고, CUDA·PyTorch·ONNXRuntime 버전을 고정한 \textbf{단일 워커 이미지}와 사내 SDK 제출 API로 학습 환경을 일원화(자동 checkpoint·중단 지점 재개 지원). \textbf{도입 단계이며 연구 코드 이관이 남은 과제} — 이관을 막는 원인을 실측으로 규명한 것이 아래 두 항목.}
        \resumeItem{\textbf{[설계 · 2026.08--현재] 기존 연구 워크플로를 실측 분석해 재현성 아키텍처를 설계.} 연구 저장소를 정량 조사해 \textbf{원격 브랜치 51개 · 출력이 저장된 노트북 셀 67\textasciitilde{}88\% · 실험 기록이 로컬 폴더에 고립}됨을 규명하고, 원인이 도구 부재가 아니라 \textbf{결과에 좌표가 없다}는 점임을 특정. 데이터셋에 \textbf{버전 · 내용 해시 · 코드 커밋 · 환경 lock · 시드}를 실행 시점에 박는 재현성 계약(동결 도구 · manifest 스키마 · 무결성 검증 CLI)을 설계·구현하고, MLflow 트래킹·모델 레지스트리(S3 아티팩트 · SSO · DataHub 미러)를 그 좌표의 저장소로 배치.}
        \resumeItem{\textbf{[기반] 학습 데이터 전달 시간을 25배 단축} — 네트워크·스토리지 재설계로 iperf3 실측 \textbf{93.7\,Mbps → 2.5\,Gbps}, 8\,TB 기준 \textbf{8.4일 → 7.8시간}. \textbf{SeaweedFS}(S3 호환) 위 Iceberg 레이크하우스에 \textbf{NVMe/HDD 티어 분리}를 적용하고, \textbf{PostgreSQL + pg\_duckdb}로 Iceberg를 ad-hoc SQL로 열어 pyiceberg 코드를 작성해야 하던 분석을 \textbf{브라우저 SQL 한 줄}로 단축.}
        \resumeItem{\textbf{[하드웨어] 문제가 코드 밖에 있을 때도 직접 해결.} 앞의 25배는 원인을 \textbf{케이블까지 내려가} 찾아낸 결과. \textbf{망가진 랜선}을 찾아 갈고 필요한 \textbf{스위치를 직접 골라 샀으며}, 사내에 \textbf{놀고 있던 RAM·NVMe를 모아 노드에 꽂아} 사양을 올림. 선과 부품부터 \textbf{Kubernetes 네트워크 부담}까지 \textbf{하드웨어든 소프트웨어든 가리지 않고} 처리.}
        \resumeItem{\textbf{[접근] 공유 비밀번호·무인증 접근을 단일 신원으로 통합.} \textbf{Dex(GitHub OIDC) + oauth2-proxy}로 사내 서비스 \textbf{23종}에 SSO를 적용하고 같은 신원을 \textbf{Kubernetes RBAC}까지 연결. GPU 서버에 직접 SSH하던 개발 방식을 \textbf{JupyterHub 브라우저 개발환경}(개인 볼륨 · egress 정책 격리 · 명령 한 줄 부트스트랩)으로 대체.}
        \resumeItem{\textbf{[판단] 실측으로 GPU 증설 결정을 전환.} 증설 검토 요청에 대해 DCGM·Prometheus로 10일간 측정한 결과 \textbf{가동률 5.3\% · 대기 작업 0건}이며 병목이 용량이 아니라 \textbf{유휴 세션이 점유한 1,112 GPU-시간}(점유 시간의 91\%)임을 규명, \textbf{구매 대신 할당 회수}로 방향을 전환.}
        \resumeItem{\textbf{[방식] AI 코딩 에이전트를 상시 투입해 1인 처리량을 팀 규모로 끌어올리되, 속도를 검증 장치로 받침.} 위 산출물을 6개월에 내면서 \textbf{SDK 단위 테스트 719건}, 태그/버전 불일치·개인키 커밋을 막는 \textbf{CI 게이트와 pre-commit 훅}, 사고 재발을 막는 \textbf{실패 기록 23건 · PRD 30건}을 함께 축적.}
    \resumeItemListEnd
    \resumeSubheading
    {Cellkey}{Jul 2025 -- Dec 2025}
    {Manager, Data Engineer}{서초구, 서울}
    \resumeItemListStart
        \resumeItem{\textbf{Apache Airflow 3.x}로 데이터 정제 파이프라인을 설계하고 로컬 환경에 배포함.}
        \resumeItem{\textbf{AWS Lambda · Batch · Step Functions}를 유지보수하고 최적화해 확장 가능한 데이터 워크플로를 구축함.}
    \resumeItemListEnd
    \resumeSubheading
    {Netmelon}{Mar 2025 -- Jun 2025}
    {Back-End Engineer}{Freelance}
    \resumeItemListStart
        \resumeItem{Android · iOS 인앱 결제 API 구현(구독 · 구독 갱신 · 구독 취소)}
        \resumeItem{Flask 웹 서버를 Docker 이미지로 말아 Google Cloud Run에 배포}
        \resumeItem{Firebase 리소스의 소유권을 확인하는 기능 구현}
        \resumeItem{LLM 기반 API 라우트 · 서비스 · 함수의 단위 테스트 작성}
        \resumeItem{Firebase Firestore · Cloud Storage 운용}
        \resumeItem{레거시 코드 리팩토링}
    \resumeItemListEnd
    \resumeSubheading
    {(주) Eazel}{Jun 2022 -- Sep 2024}
    {Team Lead, AI Labs}{중구, 서울}
    \resumeItemListStart
        \resumeItem{팀 리드로 미술 시장 ETL 파이프라인을 총괄 — 아키텍처 설계부터 배포 · 모니터링 · 성능 최적화까지 주도}
        \resumeItem{Flask/NextJS 기반 대시보드 구축}
        \resumeItem{옥션 데이터 스크래퍼 개발 및 운영}
        \resumeItem{Airflow로 데이터 수집 · 정제 · 적재를 주기 실행하는 파이프라인 구현}
        \resumeItem{AWS EMR에서 PySpark로 데이터 정제 · 적재 파이프라인 구현}
        \resumeItem{OpenAI · Facebook 임베딩 모델로 데이터 중복 제거 시스템을 개발하고, 벡터 DB를 적용해 정확도와 속도를 개선}
        \resumeItem{AWS 서비스 관리(EC2, IAM, S3, Route53, RDS, EMR)}
        \resumeItem{Github Actions로 CI/CD 구축}
        \resumeItem{LangChain · Streamlit으로 챗봇 프로토타입 구현}
        \resumeItem{챗봇용 Neo4j Cypher 쿼리 도구와 PostgreSQL 쿼리 도구 구현 주도}
        \resumeItem{LangChain과 정보 추출 프롬프트로 데이터 스크래핑 구현}
    \resumeItemListEnd

    \resumeSubheading
    {(주) Eazel}{Oct 2021 -- Jun 2022}
    {MLOps Engineer, MLOps}{중구, 서울}
    \resumeItemListStart
        \resumeItem{HuggingFace NLP 등으로 데이터 분석 → 모델링 → 검증 전 과정을 직접 수행}
        \resumeItem{인스타그램 데이터 스크래퍼 개발 및 운영}
        \resumeItem{Flask/NextJS/CytoscapeJS로 소셜 관계망 분석 대시보드를 개발해 분석 결과를 제공 — Series A 투자 유치에 기여}
        \resumeItem{DVC로 ML 라이프사이클 구현}
        \resumeItem{PostgreSQL로 데이터 관리}
    \resumeItemListEnd
    
    \resumeSubheading
    {(주) Eazel}{Jan 2019 -- Oct 2021}
    {iOS Developer, iOS Team}{강남, 서울}
    \resumeItemListStart
        \resumeItem{iOS 앱 성능과 확장성을 고려해 API 클라이언트와 연동 백엔드 로직 설계 주도}
        \resumeItem{Github Actions로 CI/CD 도입}
        \resumeItem{ReactorKit 기반 반응형 프로그래밍을 도입하고 아키텍처 설계를 주도해 앱 반응성과 유지보수성을 높임}
    \resumeItemListEnd
    
    \resumeSubheading
    {(주) Finple}{Jan 2018 -- Nov 2018}
    {Software Engineer, DevOps}{영등포, 서울}
    \resumeItemListStart
        \resumeItem{하이브리드 iOS, Android 앱 유지보수}
        \resumeItem{NodeJS 기반 데이터 스크래퍼 개발}
    \resumeItemListEnd
\resumeSubHeadingListEnd